% Part: history
% Chapter: biographies 
% Section: kurt-goedel

\documentclass[../../../include/open-logic-section]{subfiles}

\begin{document}

\olfileid{his}{bio}{god}

\olsection{کرٹ گوڈل}

\olphoto{goedel-kurt}{کرٹ گوڈل}

کرٹ گوڈل (\textsc{گَر} ڈل) 28 اپریل 1906 نوں آسٹرو ہنگیرین سلطنت دے برُن شہر (ہُن چیک جمہوریہ دا برنو) وچ جمیا۔ نِکّے کرٹیلے دی جانن دی تانگھ تے ذہانت کرکے گھر دے لوک اکثر اوہنوں «Der kleine Herr Warum» (نِکّا مسٹر کیوں) آکھدے سن۔ پرائمری سکول توں ای اوہ پڑھائی وچ اگے سی۔ صرف ریاضی وچ اوہنوں سبھ توں اُچا گریڈ نہیں سی ملدا۔ ماڑی صحت کارن گوڈل اکثر سکول توں غیر حاضر رہندا سی تے اوہنوں جسمانی تعلیم توں چھوٹ سی۔ بچپن وچ اوہنوں گٹھیا دا بخار تشخیص ہویا۔ ڈاکٹری معائنے دے اُلٹ، اوہ ساری عمر ایہ منّدا رہیا کہ ایس نال اوہدے دل اُتے پکا اثر پیا سی۔

گوڈل نے 1924 وچ ویانا یونیورسٹی وچ پڑھائی شروع کیتی تے~1929 وچ ڈاکٹریٹ دی پڑھائی پوری کیتی۔ اوہدا پہلا ارادہ طبیعیات پڑھن دا سی، پر چھیتی اوہدی دلچسپی ریاضی ول، خاص کر کے منطق ول، مُڑ گئی۔ ایہدے وچ فلسفی روڈولف کارناپ دا وی کجھ اثر سی۔ ہانس ہان دی نگرانی تھلّے لکھے اوہدے مقالے نے شناخت والی پہلے درجے دی محمولی منطق دی تکمیل دا قضیہ ثابت کیتا \citep{Godel1929}۔ بس اک سال پچھوں اوہنے اپنے سبھ توں مشہور نتیجے، نامکملیت دے پہلے تے دوجے قضیے، حاصل کیتے (ایہ \citealt{Godel1931} وچ چھپے)۔ ویانا وچ رہندیاں گوڈل ویانا سرکل نال ڈونگھا جڑیا ہویا سی۔ ایہ سائنسی سوچ والے فلسفیاں دی جماعت سی جیہدے وچ کارناپ وی سی، تے گوڈل دے نتیجیاں دا کارناپ دے کم اُتے خاص اثر پیا۔

1938 وچ گوڈل نے اڈیل نمبورْسکی نال ویاہ کیتا۔ اوہدے ماں پیو خوش نہیں سن۔ اوہ گوڈل توں چھ سال وڈی سی، اوہنوں طلاق ہو چکی سی، تے اوہ رات دے کلب وچ رقاصہ سی۔ پر سماجی دباؤ دا گوڈل اُتے اثر نہیں پیا۔ اوہ گوڈل دی وفات تک خوشی نال ویاہے رہے۔

1938 وچ نازی جرمنی دے آسٹریا نوں اپنے نال ملاؤن پچھوں گوڈل تے اڈیل امریکہ ہجرت کر گئے۔ اوتھے گوڈل نے پرنسٹن، نیو جرسی دے انسٹیٹیوٹ فار ایڈوانسڈ سٹڈی وچ عہدہ سنبھالیا۔ گھٹ میل جول تے وکھرے مزاج دے باوجود پرنسٹن وچ اوہنے ہورناں نال رَل کے بہت نتیجہ خیز کم کیتا۔ اوہنے سیٹ نظریے، فلسفے تے طبیعیات اُتے مضمون چھپوائے۔ خاص گل ایہ کہ انسٹیٹیوٹ دے اپنے ساتھی البرٹ آئن سٹائن نال اوہدی بڑی گُوڑھی دوستی ہو گئی۔

آخری سالاں وچ گوڈل دی ذہنی صحت وگڑ گئی۔ 1977 وچ اوہدی گھر والی دے ہسپتال داخل ہون نال اوہ اوہدے لئی کھانا نہیں سی پکا سکدی۔ گوڈل ساری عمر ذہنی صحت دیاں اوکڑاں نال لڑدا رہیا سی تے آخر اوہنوں شدید وہم ہو گیا۔ زہر دتے جان دے ڈر نال اوہنے کھان توں انکار کر دتا۔ اوہ 14 جنوری 1978 نوں پرنسٹن وچ بھکھ نال مر گیا۔

\begin{reading}
گوڈل دی پوری سوانح عمری لئی \citet{Dawson1997} ویکھو۔ اوہدی زندگی بارے ہور لکھتاں تے منطق تے فلسفے وچ اوہدے کم بارے مضموناں لئی \citet{Wang1990}، \citet{Baaz2011}، \citet{Takeuti2003} تے \citet{Sigmund2007} ویکھو۔

گوڈل دا ڈاکٹریٹ مقالہ اصل جرمن وچ ملدا اے \citep{Godel1929}۔ نامکملیت دے قضیاں دی اصل لکھت \citep{Godel1931} اے۔ گوڈل دیاں ساریاں چھپیاں تے ان چھپیاں لکھتاں تے چُنے ہوئے خط انگریزی وچ اوہدیاں \emph{Collected Papers} وچ ملدے نیں \citet{Godel1986,Godel1990}۔

گوڈل دے نامکملیت دے قضیاں دے تفصیلی بیان لئی \citet{Smith2013} ویکھو۔ اوہدے قضیاں دی غیر رسمی، فلسفیانہ گل بات لئی مارک لِنسن مائر دا پوڈکاسٹ ویکھو \citep{Linsenmayer2014}۔
\end{reading}

\end{document}
